\documentclass[11pt,a4paper]{article}
\usepackage[utf8]{inputenc}
\usepackage[portuguese]{babel}
\usepackage{amsmath}
\usepackage{amssymb}
\usepackage{graphicx}
\usepackage{listings}
\usepackage{xcolor}
\usepackage{float}
\usepackage{geometry}
\usepackage{fancyhdr}
\usepackage{hyperref}
\usepackage{array}
\usepackage{booktabs}

\geometry{left=1.5cm,right=1.5cm,top=1.5cm,bottom=1.5cm}
\setlength{\parskip}{0.3em}
\pagestyle{fancy}

\lhead{CC8550 --- Simulação e Teste de Software}
\chead{Aula 05}
\rhead{2º Semestre 2025}
\cfoot{Página \thepage}

\title{
    {\Large \textbf{Simulação e Teste de Software}} \\
    {\Large \textbf{Aula 05 --- Teste de Mutação (Caixa Eletrônico)}}
}

\author{
    \textbf{Aluno:} Lucas Rebouças Silva \\
    \textbf{R.A.:} 22.122.048-6 \\[0.3cm]
    \textbf{Professor:} Luciano Rossi \\
    \textbf{Centro Universitário FEI} \\
    Ciência da Computação
}

\date{2º Semestre de 2025}

\begin{document}

\maketitle

\newpage
\tableofcontents
\newpage

\section{Introdução}

Este relatório descreve o Laboratório 05 da disciplina CC8550 (Simulação e Teste de Software), cujo objetivo é aplicar \emph{teste de mutação} para avaliar a qualidade de uma suíte de testes. Foi desenvolvido um programa simples de caixa eletrônico com cinco funcionalidades e, em seguida, criados mutantes manualmente e configurada a ferramenta \texttt{mutmut} para mutação automática. Ao executar os testes no código original, todos passam; ao executá-los no código mutado, as falhas indicam que os mutantes foram ``mortos'', permitindo calcular o Mutation Score (MS) e comparar com os resultados do \texttt{mutmut} quando disponíveis.

\section{Descrição do Programa}

O sistema foi implementado no arquivo \texttt{caixa\_eletronico.py}. Uma \textbf{conta} é representada por um \emph{dataclass} com: \texttt{numero} (identificador), \texttt{saldo} (float) e \texttt{historico} (lista de strings com as operações realizadas). As cinco funcionalidades são:

\begin{enumerate}
  \item \texttt{consultar\_saldo(conta)} --- retorna o saldo atual;
  \item \texttt{depositar(conta, valor)} --- credita o valor e registra no histórico (valor deve ser positivo);
  \item \texttt{sacar(conta, valor)} --- debita o valor e registra (valor positivo e saldo suficiente);
  \item \texttt{transferir(conta\_origem, conta\_destino, valor)} --- debita da origem e credita no destino, com registro em ambas (origem e destino devem ser contas diferentes; valor positivo);
  \item \texttt{extrato(conta)} --- retorna a lista do histórico na ordem cronológica.
\end{enumerate}

Regras de negócio: depósito e saque rejeitam valor $\leq 0$; saque rejeita valor maior que o saldo; transferência rejeita mesma conta e valor $\leq 0$. O histórico de cada conta é independente.

\section{Estratégia de Testes}

Os testes foram escritos com \texttt{pytest} no arquivo \texttt{test\_caixa\_eletronico.py}. Existem \textbf{dez} casos de teste, listados na Tabela~\ref{tab:testes}, cobrindo as cinco funcionalidades e cenários de sucesso e de exceção.

\begin{table}[H]
  \centering
  \caption{Casos de teste e regras verificadas}
  \label{tab:testes}
  \begin{tabular}{lll}
    \toprule
    Teste & Funcionalidade & O que verifica \\
    \midrule
    \texttt{test\_saldo\_inicial\_zero} & consultar\_saldo & Saldo inicial 0 \\
    \texttt{test\_deposito\_aumenta\_saldo} & depositar & Depósito aumenta saldo \\
    \texttt{test\_nao\_permite\_deposito\_negativo\_ou\_zero} & depositar & Rejeita 0 e valor negativo \\
    \texttt{test\_saque\_diminui\_saldo} & sacar & Saque debita corretamente \\
    \texttt{test\_nao\_permite\_saque\_maior\_que\_saldo} & sacar & Rejeita saque acima do saldo \\
    \texttt{test\_pode\_sacar\_todo\_saldo} & sacar & Permite sacar 100\% do saldo \\
    \texttt{test\_transferencia\_move\_saldo\_correto} & transferir & Valores corretos origem/destino \\
    \texttt{test\_nao\_permite\_transferencia\_para\_mesma\_conta} & transferir & Rejeita mesma conta \\
    \texttt{test\_extrato\_registra\_operacoes\_na\_ordem} & extrato & Ordem cronológica do histórico \\
    \texttt{test\_contas\_tem\_historicos\_independentes} & extrato & Históricos não compartilhados \\
    \bottomrule
  \end{tabular}
\end{table}

Ao rodar \texttt{pytest} no código original, os dez testes passam.

\section{Mutantes Manuais}

\subsection*{O que são mutantes mortos e sobreviventes?}

Um \textbf{mutante morto} é uma versão do programa em que foi introduzida uma alteração (mutação) e, ao rodar a suíte de testes, \emph{pelo menos um teste falha}. Isso indica que os testes detectaram a diferença entre o código original e o mutado: o defeito simulado foi revelado e o mutante é considerado ``morto'' pela suíte. Já um \textbf{mutante sobrevivente} é aquele em que todos os testes continuam passando mesmo com o código alterado; nesse caso, a suíte não conseguiu distinguir o mutante do programa original, o que sugere lacuna na cobertura ou na assertividade dos testes. O objetivo é ter uma suíte que mate o máximo possível de mutantes (alto Mutation Score).

Foi criada uma cópia do código em \texttt{caixa\_eletronico\_mutante.py} com \textbf{quatro} mutantes manuais:

\begin{itemize}
  \item \textbf{Mutante 1 (depositar):} condição alterada de \texttt{valor <= 0} para \texttt{valor < 0}, permitindo depósito de R\$ 0,00. Morto por \texttt{test\_nao\_permite\_deposito\_negativo\_ou\_zero}.
  \item \textbf{Mutante 2 (sacar):} condição \texttt{valor > conta.saldo} trocada por \texttt{valor >= conta.saldo}, impedindo sacar todo o saldo. Morto por \texttt{test\_pode\_sacar\_todo\_saldo}.
  \item \textbf{Mutante 3 (transferir):} ordem invertida --- \texttt{sacar(conta\_destino, valor)} e \texttt{depositar(conta\_origem, valor)}. Morto por \texttt{test\_transferencia\_move\_saldo\_correto}.
  \item \textbf{Mutante 4 (extrato):} \texttt{return list(conta.historico)} trocado por \texttt{return list(reversed(conta.historico))}. Morto por \texttt{test\_extrato\_registra\_operacoes\_na\_ordem}.
\end{itemize}

Ao executar a mesma suíte de testes no código mutante (importando \texttt{caixa\_eletronico\_mutante}), vários testes falham; cada mutante é detectado por pelo menos um teste. Fórmula do Mutation Score:
\[
MS = \frac{\text{mutantes mortos}}{\text{mutantes totais} - \text{mutantes equivalentes}} \times 100\%.
\]

\begin{table}[H]
  \centering
  \caption{Resultados dos mutantes manuais}
  \label{tab:mutantes-manuais}
  \begin{tabular}{lcc}
    \toprule
    & Quantidade & \\
    \midrule
    Mutantes totais                & 4 & \\
    Mutantes equivalentes         & 0 & \\
    Mutantes mortos               & 4 & \\
    Mutantes sobreviventes        & 0 & \\
    \midrule
    \textbf{Mutation Score (MS)}  & $4/4 = 100\%$ & \\
    \bottomrule
  \end{tabular}
\end{table}

\section{Mutação Automática com \texttt{mutmut}}

A ferramenta \texttt{mutmut} foi configurada no \texttt{pyproject.toml}: mutação em \texttt{caixa\_eletronico.py}, diretório de testes no próprio projeto e runner \texttt{python -m pytest}. Os comandos para execução são \texttt{mutmut run} (gera e executa os mutantes) e \texttt{mutmut results} (exibe o resumo). No ambiente Windows o \texttt{mutmut} não é suportado nativamente; a documentação oficial recomenda o uso de WSL (Windows Subsystem for Linux) ou um ambiente Linux. Ao rodar em WSL/Linux, preencher a tabela abaixo com os números exibidos por \texttt{mutmut results} e calcular o MS pela mesma fórmula.

\begin{table}[H]
  \centering
  \caption{Resultados do \texttt{mutmut} (preencher após executar em WSL/Linux)}
  \label{tab:mutmut}
  \begin{tabular}{lcc}
    \toprule
    & Quantidade & \\
    \midrule
    Mutantes totais                & (preencher) & \\
    Mutantes equivalentes         & (preencher) & \\
    Mutantes mortos                & (preencher) & \\
    Mutantes sobreviventes        & (preencher) & \\
    \midrule
    \textbf{Mutation Score (MS)}   & (preencher) & \\
    \bottomrule
  \end{tabular}
\end{table}

\section{Discussão e Conclusão}

Nos mutantes manuais, a suíte de testes matou os quatro mutantes introduzidos, resultando em Mutation Score de 100\% para esse conjunto. Isso indica que os testes cobrem bem as regras alteradas: rejeição de depósito zero/negativo, possibilidade de sacar todo o saldo, direção correta da transferência e ordem cronológica do extrato. O \texttt{mutmut} gera mutantes automaticamente (troca de operadores, constantes, etc.) em maior quantidade; ao executá-lo em WSL, é possível obter um MS global e comparar com o resultado manual. Mutantes sobreviventes do \texttt{mutmut} sugerem pontos em que a suíte poderia ser reforçada com novos testes ou asserções mais específicas. Em resumo, o teste de mutação complementa a cobertura estrutural ao avaliar se os testes são capazes de detectar defeitos simulados, e neste laboratório a abordagem manual mostrou-se eficaz para os quatro cenários escolhidos.

\end{document}
